\begin{abstract}
% 150–250 words
\glspl{ccTLD} represent national digital identities and are often used to provide local content to the population. Governmental domains under \glspl{ccTLD} hold inherent trust from citizens, and government-to-citizen communications are crucial. As a result, phishing attacks that abuse governmental domains can severely undermine public trust and the credibility of government communications.
The adoption of secure email mechanisms such as \gls{DMARC} is vital to combating phishing campaigns.
While \gls{DMARC} deployment has been studied broadly, little is known about its adoption and reporting practices from a geographical and \gls{ccTLD}-specific perspective.
In this paper, we analyze \gls{DMARC} policy data of domain names under 301 ccTLDs, collected at regular intervals over $\approx$1.5 years, and enriched with government domain sources.
We assess adoption trends, policy strength, and the global flow of \gls{DMARC} reports, highlighting dependencies on foreign service providers.
Our results show that while European countries lead in \gls{DMARC} adoption, domains using a stringent policy remain limited at 34.5\% (5.2M).
Within government domains, we find that \gls{DMARC} reporting data is configured to cross national borders and, in cases, relies on third-party processors -- often U.S.-based -- raising digital sovereignty concerns. 21.8\% (5.9k) non-U.S. governmental domains send \gls{DMARC} reports to U.S.-based mail boxes.
We conclude with targeted recommendations for different stakeholders, including national authorities, registries, and email service providers, to strengthen national email security.

% IEEE style
\begin{IEEEkeywords}
Email anti-spoofing, DMARC, digital sovereignty
\end{IEEEkeywords}
\end{abstract}
